% ch23-hangfire-jobs.tex — الفصل: Hangfire والمهامّ الخلفيّة (Background Jobs)
\section{\en{Hangfire} والمهامّ الخلفيّة (\en{Background Jobs})}

في الفصل ده هنفهم حاجة بسيطة بس مهمّة جدًّا لأي نظام حقيقيّ: مش كل شغل لازم يتعمل
«وإحنا ماسكين المستخدم منتظر». في شغل تقيل أو مؤجَّل الأحسن نعمله في الخلفيّة، بعيدًا عن
الطلب اللي جاي من المتصفّح. هنعرف ليه نحتاج ده، وإيه هي أداة \en{Hangfire} اللي بنستعملها،
وإيه المهامّ الحقيقيّة اللي شغّالة في \en{ScholarPath}، وإزاي المهمّة الواحدة بتمشي خطوة
بخطوة من لحظة ما نطلبها لحد ما تخلص. القارئ يفترض إنه فاهم من قبل فكرة الطبقات
(\en{layers})، والـ\en{ports}، والـ\en{idempotency} — وهنعتمد عليهم من غير ما نعيد شرحهم
بعمق.

% ═══════════════════════════════════════════════════════════════════════
\subsection{ليه نحتاج مهامّ خلفيّة أصلًا؟}

خلّينا نتخيّل إن طالب ضغط زرّ «قدِّم على المنحة». في الخلفيّة لازم يحصل: نحفظ الطلب،
ونبعت له إيميل تأكيد، ويمكن نحدّث بعض الإحصاءات. لو عملنا كل ده جوّه الطلب الواحد،
الطالب هيفضل يبصّ في شاشة بتلفّ لحد ما الإيميل يتبعت والإحصاء يتحسب. وده بطيء وبيتعطّل
لو خدمة الإيميل زعلت.

الفكرة إن الشغل نوعين. في شغل «متزامن» (\en{synchronous}) لازم يتعمل جوّه الطلب لأن الردّ
يعتمد عليه — زي التحقّق من الصلاحيّة وحفظ الطلب نفسه. وفي شغل «خلفيّ» (\en{background})
ممكن يتأجّل لأن المستخدم مش محتاج ينتظره — زي إرسال الإيميل، أو معالجة تسجيل فيديو
الجلسة، أو تجهيز تقرير. النوع التاني بنرميه لطابور خلفيّ ونردّ على المستخدم فورًا.

\begin{idea}[الفكرة: افصل الثقيل عن الطلب]
القاعدة البسيطة: لو المستخدم مش محتاج نتيجة العمليّة عشان يكمّل، متخلّيهوش يستنّاها.
احفظ اللي لازم يتحفظ فورًا، وارمي الباقي (الإيميل، الفيديو، التقرير) للخلفيّة، وردّ عليه
بسرعة.
\end{idea}

وفي كمان نوع تالت مهمّ: شغل «مجدوَل» (\en{scheduled}) بيشتغل من نفسه على مواعيد، من غير
ما حد يضغط أي زرّ. زي تذكيرات ليليّة باقتراب مواعيد المنح، أو كنس دوريّ للجلسات اللي حصل
فيها غياب (\en{no-show}). ده شغل مالوش «طلب مستخدم» أصلًا — محتاج حاجة تصحّيه بنفسها كل
فترة.

\begin{center}\begin{LTR}
\begin{tikzpicture}[node distance=0.9cm and 1.3cm]
  \node[actor] (user) {User};
  \node[pbox, right=of user] (api) {API handler};
  \node[tbox, right=of api] (resp) {fast 202 reply};
  \node[sbox, below=1.2cm of api] (enq) {enqueue};
  \node[dbbox, right=of enq] (store) {job storage};
  \node[pbox, right=of store] (worker) {worker: email / video};

  \draw[arr] (user) -- node[lbl,above]{request} (api);
  \draw[arr] (api) -- node[lbl,above]{fast} (resp);
  \draw[arr] (resp.south) to[bend left=25] node[lbl,right]{back} (user.east);
  \draw[arrd] (api) -- node[lbl,right]{offload heavy} (enq);
  \draw[arr] (enq) -- (store);
  \draw[arr] (store) -- (worker);
\end{tikzpicture}
\end{LTR}\end{center}

الرسمة فوق بتقول القصّة كلها: الطلب بيدخل، الـ\en{API} بيعمل اللازم بسرعة ويردّ (السهم
العلويّ للمستخدم)، وفي نفس الوقت بيرمي الشغل التقيل لطابور خلفيّ (السهم المتقطّع لأسفل)،
والعامل (\en{worker}) بيلتقطه ويشتغل عليه بعيدًا. المستخدم مستنّاش الجزء التقيل خالص.

\begin{spotlight}[ركّز هنا لإبهار اللجنة: ردّ سريع للمستخدم]
أهمّ رسالة تقولها للّجنة: «إحنا بنردّ على المستخدم بسرعة لأن الشغل التقيل بيتأجّل
للخلفيّة». تقديم طلب منحة بيرجع للطالب في أجزاء من الثانية، بينما إرسال الإيميل ومعالجة
الفيديو وتجهيز التقارير بتحصل بعيدًا عن الطلب. النتيجة: واجهة سريعة لا تتعطّل لو خدمة
خارجيّة (زي بريد أو \en{Stripe}) بطُؤت لحظيًّا.
\end{spotlight}

% ═══════════════════════════════════════════════════════════════════════
\subsection{يعني إيه Hangfire؟}

\en{Hangfire} مكتبة لإدارة المهامّ الخلفيّة في \en{.NET}. بدل ما نبني طابور بإيدينا، هي
بتدّينا أربع حاجات جاهزة:

\begin{itemize}
  \item طابور مدعوم بتخزين دائم — المهامّ مش بتتخزّن في الذاكرة (اللي بتضيع لو
  السيرفر وقع)، لكن في جداول جوّه قاعدة البيانات (\en{SQL Server} عندنا). يعني لو السيرفر
  اتعمل له \en{restart}، المهامّ لسه موجودة وهتكمّل.
  \item عمّال (\en{workers}) — خيوط بتفضل تسأل الطابور: «في شغل؟» وتلتقطه وتنفّذه.
  \item جدولة — تقدر تجدول مهمّة تشتغل بعد شويّة، أو تتكرّر بجدول \en{cron}.
  \item لوحة تحكّم (\en{dashboard}) — صفحة بتوريك المهامّ الشغّالة والفاشلة
  والمعادة، ومحميّة عندنا بصلاحيّة \en{Admin} بس.
\end{itemize}

وفي تلات طرق أساسيّة نطلب بيها مهمّة:

\begin{center}
\begin{tabular}{L{3.4cm} L{11.2cm}}
\toprule
النوع & المعنى \\
\midrule
\en{enqueue} (\en{fire-and-forget}) & نفّذها دلوقتي في الخلفيّة، وانسَها — زي إرسال إيميل بعد حفظ الطلب. \\
\en{schedule} & نفّذها مرّة واحدة بعد وقت محدّد (مثلًا بعد ساعة). \\
\en{recurring} & كرّرها على جدول \en{cron} — زي كل 15 دقيقة أو كل ليلة. \\
\bottomrule
\end{tabular}
\end{center}

وأهمّ ميزة: لو المهمّة فشلت (استثناء)، \en{Hangfire} بيعيد محاولتها تلقائيًّا كذا مرّة
بفواصل متزايدة. وده بيفتح سؤال مهمّ: لو المهمّة اتنفّذت نصّها وفشلت وأُعيدت، هتتعمل مرّتين؟
عشان كده لازم كل مهمّة تكون \en{idempotent} — يعني تنفيذها مرّتين بنفس المدخلات يدّي نفس
النتيجة من غير ضرر (لا إيميل مكرّر ولا دفعة مزدوجة). هنرجع للنقطة دي بالتفصيل في آخر الفصل.

\begin{keyp}[نقطة مهمة: التخزين هو اللي بيخلّي الطابور ينجو من السقوط]
لأن \en{Hangfire} بيكتب المهامّ في قاعدة البيانات مش في الذاكرة، أي مهمّة اتحطّت في الطابور
بتفضل موجودة حتى لو السيرفر وقع أو اتعمل له \en{deploy}. ده الفرق بين طابور «هشّ» في
الذاكرة وطابور «دائم» يعتمد عليه.
\end{keyp}

الكود ده بيبيّن إزاي بنركّب \en{Hangfire} على تخزين \en{SQL Server}. لاحظ إن كل ده متعلّق
بمفتاح إعداد — لو الميزة مقفولة، مفيش تخزين ولا عمّال أصلًا.

\begin{english}
\begin{lstlisting}[language={[Sharp]C}]
builder.Services.AddHangfire(cfg =>
{
    cfg.UseSimpleAssemblyNameTypeSerializer()
       .UseRecommendedSerializerSettings();

    if (!string.IsNullOrWhiteSpace(connectionString))
    {
        cfg.UseSqlServerStorage(connectionString, new SqlServerStorageOptions
        {
            CommandBatchMaxTimeout = TimeSpan.FromMinutes(5),
            QueuePollInterval = TimeSpan.FromSeconds(30),
        });
    }
    // ...
});

builder.Services.AddHangfireServer();   // start the workers
\end{lstlisting}
\end{english}

شرح سريع لأهمّ الأسطر: \code{UseSqlServerStorage} بيقول لـ\en{Hangfire} يخزّن جداوله في
نفس قاعدة البيانات (نفس \code{connectionString}). و\code{QueuePollInterval} هو كل قدّ إيه
العامل يسأل الطابور عن شغل جديد (كل 30 ثانية). و\code{AddHangfireServer} هو اللي بيشغّل
العمّال اللي بيلتقطوا المهامّ فعلًا. لو التخزين مش متاح، بيرجع لتخزين ذاكرة مؤقّت
(\en{memory storage}) للتطوير.

% ═══════════════════════════════════════════════════════════════════════
\subsection{المهامّ الحقيقيّة في ScholarPath}

كل المهامّ المتكرّرة متحكَّم فيها بمفتاح واحد: \code{Hangfire:Enabled}. لو \code{false}
(القيمة الافتراضيّة في التطوير)، مفيش عمّال بيتشغّلوا ولا مهامّ بتتجدول خالص — ده بيخلّي
التشغيل المحلّي خفيف. في الإنتاج بنشغّله. الجدول ده بيلخّص أهمّ المهامّ الفعليّة المسجَّلة،
وجدولها، وبتعمل إيه:

\begin{center}
\begin{tabular}{L{3.7cm} L{2.7cm} L{7.6cm}}
\toprule
المهمّة & الجدول & بتعمل إيه \\
\midrule
\en{deadline-reminders} & يوميًّا 09:00 & تذكّر الطلاب باقتراب موعد المنح المحفوظة أو المتقدَّم لها. \\
\en{draft-reminders} & يوميًّا 10:00 & تنكز مَن عنده مسودّة طلب لسه ما اتبعتتش. \\
\en{scholarship-auto-close} & كل ساعة & تقفل المنح تلقائيًّا بعد ما موعدها يعدّي. \\
\en{meeting-no-show-sweep} & كل 15 دقيقة & ترصد غياب أحد طرفَي الجلسة وتفتح \en{NoShowReport}. \\
\en{booking-reminders} & كل 15 دقيقة & تبعت تذكير قبل الجلسة بـ 24 ساعة وبساعة. \\
\en{stripe-payouts} & يوميًّا 02:00 & تصرف دفعات \en{Stripe} للمستشارين على دفعة مجمّعة. \\
\en{data-export-sweep} & كل ساعة & تنفّذ طلبات تصدير بيانات \en{GDPR}. \\
\en{data-delete-sweep} & يوميًّا 03:00 & تنفّذ طلبات حذف بيانات \en{GDPR}. \\
\en{redaction-audit-sampling} & شهريًّا (اليوم 1، 02:00) & تسحب عيّنة تدقيق لتنقية بيانات \en{PII}. \\
\en{integrity-check} & يوميًّا 04:00 & فحص سلامة يكشف الصفوف اليتيمة وغير المتّسقة. \\
\bottomrule
\end{tabular}
\end{center}

\begin{note}
فيه شغل خلفيّ مش متكرّر بجدول، لكنه بيتحطّ في الطابور «عند الطلب» — أشهر مثال إعادة بناء
قاعدة المعرفة للذكاء الاصطناعيّ (\en{RebuildKnowledgeBase}). ده بيتفعّل لمّا الأدمن يضغط
الزرّ (أو مرّة عند إقلاع النظام) بدل ما يتكرّر كل ليلة، لأنه شغل تقيل مش محتاج يتعمل كل يوم.
\end{note}

الكود ده هو تسجيل المهامّ المتكرّرة نفسه. لاحظ إن كله جوّه شرط \code{Hangfire:Enabled}،
وإن كل مهمّة ليها اسم ثابت وجدول \en{cron}:

\begin{english}
\begin{lstlisting}[language={[Sharp]C}]
if (hangfireOpts.Enabled)
{
    var recurring = app.Services.GetRequiredService<IRecurringJobManager>();

    recurring.AddOrUpdate<IMeetingNoShowSweepJob>(
        "meeting-no-show-sweep", j => j.RunAsync(CancellationToken.None), "*/15 * * * *");
    recurring.AddOrUpdate<IStripePayoutJob>(
        "stripe-payouts", j => j.RunAsync(CancellationToken.None), Cron.Daily(2));   // 02:00 UTC
    recurring.AddOrUpdate<IScholarshipAutoCloseJob>(
        "scholarship-auto-close", j => j.RunAsync(CancellationToken.None), Cron.Hourly);
    recurring.AddOrUpdate<IDeadlineReminderJob>(
        "deadline-reminders", j => j.RunAsync(CancellationToken.None), Cron.Daily(9));
    // ... the rest of the jobs follow the same pattern
}
\end{lstlisting}
\end{english}

شرح: \code{AddOrUpdate} بيسجّل مهمّة متكرّرة باسم ثابت — لو موجودة بيحدّثها بدل ما يعمل
نسخة تانية. الوسيط التالت هو الجدول: \code{"*/15 * * * *"} تعبير \en{cron} معناه «كل 15
دقيقة»، و\code{Cron.Daily(2)} معناها الساعة 02:00 بتوقيت \en{UTC}. لاحظ إن كل مهمّة بتُسجَّل
عبر \en{interface} زي \code{IStripePayoutJob} — وده مهمّ هنرجع له في آخر الفصل. الرسمة دي
بتوزّع المهامّ على «خطّ زمنيّ» حسب تكرارها:

\begin{center}\begin{LTR}
\begin{tikzpicture}[node distance=0.55cm and 0.5cm]
  \node[gbox] (t1) {every 15 min};
  \node[wbox, right=of t1] (r1) {no-show sweep, booking reminders, booking completion, session expiry};
  \node[gbox, below=of t1] (t2) {hourly};
  \node[wbox, right=of t2] (r2) {scholarship auto-close, GDPR export sweep};
  \node[gbox, below=of t2] (t3) {nightly (02:00 - 10:00)};
  \node[wbox, right=of t3] (r3) {Stripe payouts, GDPR delete, integrity check, deadline and draft reminders};
  \node[gbox, below=of t3] (t4) {monthly (1st, 02:00)};
  \node[wbox, right=of t4] (r4) {PII redaction audit sampling};

  \draw[arr] (t1) -- (t2);
  \draw[arr] (t2) -- (t3);
  \draw[arr] (t3) -- (t4);
\end{tikzpicture}
\end{LTR}\end{center}

الأعمدة اليسرى (بنفسجيّ) هي الإيقاع: من الأسرع (كل 15 دقيقة) للأبطأ (شهريًّا)، والسهم النازل
بيمثّل مرور الزمن. الأشرطة الرماديّة الطويلة بتقول المهامّ اللي بتشتغل في كل إيقاع.

% ═══════════════════════════════════════════════════════════════════════
\subsection{كيف تشتغل مهمّة خطوة بخطوة}

خلّينا نتتبّع رحلة مهمّة واحدة من أوّلها لآخرها. الرحلة بسيطة: حاجة بتحطّ المهمّة في الطابور
(\en{enqueue})، التخزين بيسجّلها، عامل بيلتقطها، بينفّذها، وبعدين يا إمّا تنجح يا إمّا تفشل
وتُعاد تلقائيًّا.

\begin{center}\begin{LTR}
\begin{tikzpicture}[node distance=0.8cm and 1.05cm]
  \node[sbox] (enq) {enqueue};
  \node[dbbox, right=of enq] (q) {storage / queue};
  \node[pbox, right=of q] (w) {worker executes};
  \node[dbox, right=of w] (dec) {success?};
  \node[tbox, right=of dec] (done) {done};
  \node[fbox, below=1.1cm of dec] (retry) {retry (back-off)};

  \draw[arr] (enq) -- (q);
  \draw[arr] (q) -- (w);
  \draw[arr] (w) -- (dec);
  \draw[arr] (dec) -- node[lbl,above]{yes} (done);
  \draw[arr] (dec) -- node[lbl,right]{no} (retry);
  \draw[arrd] (retry) -| node[lbl,below]{re-queued} (q);
\end{tikzpicture}
\end{LTR}\end{center}

المهمّ في الرسمة هو الفرع النازل: لو فشلت المهمّة، مبتضيعش — بترجع للطابور (السهم المتقطّع)
عشان تتحاول تاني بعد فترة انتظار متزايدة. ودي بالظبط النقطة اللي بتخلّي الـ\en{idempotency}
شرطًا مش رفاهيّة.

الكود ده مقتطف من مهمّة تذكيرات المواعيد الحقيقيّة (\en{DeadlineReminderJob}). شوف إزاي
بتبعت التذكير — وركّز في السطر بتاع \code{idempotencyKey}:

\begin{english}
\begin{lstlisting}[language={[Sharp]C}]
foreach (var recipientId in recipients)
{
    await notifications.DispatchAsync(
        recipientId,
        NotificationType.ApplicationDeadlineApproaching,
        new NotificationParams
        {
            TitleEn = scholarship.TitleEn,
            TitleAr = scholarship.TitleAr,
            Count = daysLeft,
        },
        deepLink: $"/student/scholarships/{scholarship.Id}",
        idempotencyKey: $"deadline-reminder:{scholarship.Id:N}:{recipientId:N}:{deadlineStamp}",
        ct).ConfigureAwait(false);
}
\end{lstlisting}
\end{english}

شرح: المهمّة بتلفّ على كل مستلم وتبعت له تذكير عبر \code{DispatchAsync}. السطر الأهمّ هو
\code{idempotencyKey}: مفتاح مبنيّ من رقم المنحة، ورقم المستلم، وتاريخ الموعد
(\code{deadlineStamp}). يعني نفس التذكير لنفس الشخص لنفس الموعد ليه بصمة واحدة ثابتة. لو
المهمّة اليوميّة اشتغلت مرّتين بالغلط، أو أُعيدت بعد فشل، طبقة الإشعارات هتشوف إن البصمة دي
اتبعتت قبل كده فمش هتبعت تاني. تذكير واحد بالظبط لكل موعد — لا تكرار.

% ═══════════════════════════════════════════════════════════════════════
\subsection{Idempotency وإعادة المحاولة: لا صرف مزدوج}

دلوقتي نجمع الخيطين: إعادة المحاولة التلقائيّة نعمة، لكنها بتفتح خطر التكرار. الحلّ إن كل
مهمّة تتصمّم بحيث لو اتنفّذت مرّتين متعملش ضرر. في التذكيرات عملنا ده بمفتاح
\code{idempotencyKey}. في صرف الفلوس الموضوع أخطر بكتير — لو صرفنا لمستشار مرّتين، دي فلوس
حقيقيّة ضاعت.

عشان كده مهمّة صرف \en{Stripe} بتتحمي بطبقتين. الأولى قفل بيمنع تشغيلين متوازيين لنفس
المهمّة في نفس اللحظة. شوف إزاي القفل متحطّ على الـ\en{interface} نفسه:

\begin{english}
\begin{lstlisting}[language={[Sharp]C}]
public interface IStripePayoutJob
{
    // The lock must live HERE on the interface method — Hangfire reads
    // job filters off the registered (interface) method, not the impl.
    [DisableConcurrentExecution(timeoutInSeconds: 30 * 60)]
    Task RunAsync(CancellationToken ct);
}
\end{lstlisting}
\end{english}

شرح: \code{DisableConcurrentExecution} بيقول لـ\en{Hangfire} «مسموح بتشغيل واحد بس لهذه
المهمّة في المرّة». الطبقة التانية (في التنفيذ) هي إن كل دفعة بتتعلّم عليها علامة قبل الصرف
باستخدام رمز تزامن (\en{RowVersion}) على صفّ الدفعة — يعني حتى لو اتنين خادمين حاولوا، واحد
بس هينجح في «حجز» الدفعة والصرف. الاتنين مع بعض: لا صرف مزدوج مهما تكرّرت المهمّة.

\begin{spotlight}[ركّز هنا لإبهار اللجنة: \en{idempotency} + إعادة المحاولة = صفر تكرار]
هنا بيت القصيد اللي بيفرّق نظام هاوي عن نظام جادّ. لأن المهامّ بتتعاد تلقائيًّا عند الفشل،
صمّمنا كل مهمّة تكون \en{idempotent}: التذكيرات ببصمة \code{idempotencyKey} فمفيش إشعار
مكرّر، والدفعات بقفل \code{DisableConcurrentExecution} + رمز التزامن على الصفّ فمفيش صرف
مزدوج ولا فلوس ضايعة. الرسالة للّجنة: «إعادة المحاولة تحمينا من الفشل، والـ\en{idempotency}
تحمينا من إعادة المحاولة نفسها».
\end{spotlight}

\begin{secnote}
لوحة تحكّم \en{Hangfire} على المسار \code{/hangfire} مش مفتوحة للكل — محميّة بفلتر صلاحيّة
يسمح لـ\en{Admin} بس، وبتتفعّل بمفتاح منفصل \code{DashboardEnabled}. يعني حتى الأداة اللي
بتوريك قلب النظام الخلفيّ نفسها متقفلة خلف الصلاحيّات.
\end{secnote}

\begin{qa}[سؤال متوقّع: طب لو نفس المهمّة اليوميّة اشتغلت مرّتين؟]
مفيش مشكلة. المهامّ عندنا مصمّمة على مبدأ «آمنة عند التكرار». التذكيرات هتتصفّى بالبصمة،
والدفعات هيتصدّها القفل ورمز التزامن، والإغلاق التلقائيّ للمنح بيعتمد على الحالة الحاليّة
(\en{Open} فقط) فتشغيله تاني على منحة اتقفلت خلاص مبيعملش حاجة. ده بالظبط تعريف الـ
\en{idempotency}.
\end{qa}

\begin{why}
ليه نتكبّد كل ده بدل ما نشغّل الشغل جوّه الطلب؟ لأن الفصل ده بيدّينا تلات مكاسب مع بعض:
واجهة سريعة (المستخدم مستنّاش الثقيل)، وصمود (المهمّة في التخزين تنجو من سقوط السيرفر
وتُعاد لو فشلت)، وأتمتة (شغل بيحصل من نفسه من غير تدخّل بشريّ زي إغلاق المنح والتذكيرات).
\end{why}

% ═══════════════════════════════════════════════════════════════════════
\begin{summary}[الخلاصة في سطور]
\begin{itemize}
  \item الشغل التقيل أو المؤجَّل مايصحّش يعطّل ردّ المستخدم — بنرميه للخلفيّة ونردّ فورًا.
  \item \en{Hangfire} بيدّينا طابورًا مدعومًا بتخزين دائم في \en{SQL Server}، وعمّالًا
  يلتقطوا المهامّ، وجدولة، ولوحة تحكّم محميّة بـ\en{Admin}.
  \item تلات طرق للطلب: \en{enqueue} فوري، و\en{schedule} مؤجَّل، و\en{recurring} بجدول
  \en{cron}. وكله متحكَّم فيه بمفتاح \code{Hangfire:Enabled}.
  \item عندنا مهامّ حقيقيّة: تذكيرات المواعيد والمسودّات، الإغلاق التلقائيّ للمنح، كنس
  الـ\en{no-show}، صرف \en{Stripe}، تصدير وحذف \en{GDPR}، عيّنات تدقيق تنقية \en{PII}،
  وفحوص السلامة.
  \item دورة حياة المهمّة: \en{enqueue} $\rightarrow$ تخزين $\rightarrow$ عامل يلتقط
  $\rightarrow$ تنفيذ $\rightarrow$ نجاح أو إعادة محاولة تلقائيّة.
  \item لأن الفشل بيسبّب إعادة محاولة، كل مهمّة لازم تكون \en{idempotent}: بصمة
  \code{idempotencyKey} للإشعارات، وقفل + رمز تزامن للدفعات — فلا إشعار مكرّر ولا صرف
  مزدوج.
\end{itemize}
\end{summary}
